\documentclass[11pt]{article}

\usepackage[margin=1in]{geometry}
\usepackage{amsmath,amssymb}
\usepackage{bm}
\usepackage{booktabs}
\usepackage{enumitem}
\usepackage{hyperref}

\title{A Hierarchical Gene--Isoform Model for Sparse Single-Cell Data}
\author{}
\date{}

\begin{document}
\maketitle

\section*{Purpose}

The existing genome-wide factorization models normalized equivalence-class
(EC) counts directly. It recovers cell-type information, but a matched
evaluation shows a substantial gap from a standard gene-expression analysis.
Increasing the factor rank does not close this gap. Adding an auxiliary
gene-reconstruction loss helps, indicating that the EC objective spends too
much capacity on transcript ambiguity and too little on marker-gene
abundance.

This note proposes a hierarchical count model that separates gene abundance
from conditional isoform allocation. The model is designed for cells that do
not contain enough molecules for independent transcript-level EM. Statistical
strength instead comes from pooled transcript baselines, low-dimensional cell
states, and transcript loadings shared across all cells.

\section*{Notation}

Let \(m=1,\ldots,M\) index cells, \(g=1,\ldots,G\) genes, and
\(t=1,\ldots,T\) transcripts. Let \(g(t)\) be the gene containing transcript
\(t\), and let \(\mathcal T_g=\{t:g(t)=g\}\). Primer or library type is indexed
by \(p\). For primer \(p\), let
\[
  A_p\in\mathbb R_+^{K_p\times T}
\]
be the fixed EC-by-transcript observation matrix. It includes alignment
compatibility, Salmon conditional weights and positional-bias correction, and
the primer-specific transcript sampling exposure. In particular, the primary
paired model uses no additional transcript-length exposure for oligo(dT) and
relative effective-length exposure for random-hexamer priming.

Let \(n_{kmp}\) be the observed UMI count for EC \(k\), cell \(m\), and primer
\(p\), and let \(N_{mp}=\sum_k n_{kmp}\).

\section*{Hierarchical Abundance}

Transcript abundance is decomposed as
\[
  \theta_{tm}=a_{g(t)m}\pi_{tm},
  \qquad
  \sum_{t\in\mathcal T_g}\pi_{tm}=1.
\]
Here \(a_{gm}>0\) is the abundance of gene \(g\) in cell \(m\), while
\(\pi_{tm}\) is the conditional fraction of gene \(g(t)\)'s molecules assigned
to transcript \(t\).

This within-gene normalization is not a global simplex constraint on
\(\Theta\). It is needed to identify gene abundance separately from isoform
scale. The gene abundances may remain unnormalized because the observation
likelihood below normalizes predicted EC mass within each cell and primer.

\subsection*{Gene abundance factors}

Let \(\bm z_m\in\mathbb R^{r_g}\) be the gene-expression state of cell \(m\).
Model gene abundance as
\[
  \log a_{gm}=\alpha_g+\bm u_g'\bm z_m,
\]
where \(\alpha_g\) is a pooled gene intercept and \(\bm u_g\) is a gene loading.
The cell state \(\bm z_m\) is the natural representation for cell annotation,
visualization, and comparison with standard gene-expression analyses.

An exponential link is convenient for a count model. A softplus link can be
used if bounded gradients are more stable:
\[
  a_{gm}=\operatorname{softplus}(\alpha_g+\bm u_g'\bm z_m).
\]

\subsection*{Conditional isoform programs}

For gene \(g\), let
\[
  P_g\in\mathbb R_+^{T_g\times C_g},
  \qquad
  \bm 1'P_g=\bm 1',
\]
where \(T_g=|\mathcal T_g|\). Each column of \(P_g\) is a conditional isoform
usage program, and \(C_g\) is small, typically one to three. Let
\(\bm s_{gm}\in\mathbb R_+^{C_g}\) be the cell-specific program weights:
\[
  \lambda_{gmc}
  =
  \kappa_{gc}+\bm v_{gc}'\bm z_m,
  \qquad
  s_{gmc}
  =
  \frac{\exp(\lambda_{gmc})}
       {\sum_{d=1}^{C_g}\exp(\lambda_{gmd})}.
\]
Conditional transcript usage is
\[
  \bm\pi_{gm}=P_g\bm s_{gm}.
\]
Thus the same robust cell state \(\bm z_m\) controls gene abundance through
\(\bm u_g\) and isoform-program weights through \(\bm v_{gc}\). A cell does
not receive an unrestricted transcript vector.

If held-out molecules support additional isoform structure, introduce a
smaller isoform-specific state \(\bm h_m\in\mathbb R^{r_i}\):
\[
  \lambda_{gmc}
  =
  \kappa_{gc}+\bm v_{gc}'\bm z_m+\bm r_{gc}'\bm h_m.
\]
The \(\bm r_{gc}'\bm h_m\) term captures program changes that are not
predictable from gene expression. The nested model \(r_i=0\) is the default,
and \(r_i>0\) is accepted only by label-blind molecule-count
cross-validation.

\subsection*{Restriction relative to a flat model}

The decomposition \(\theta_{tm}=a_{g(t)m}\pi_{tm}\) alone is only a
reparameterization. Given any positive \(\bm\theta_m\), its gene totals and
within-gene proportions recover \(a_{gm}\) and \(\bm\pi_{gm}\). A multinomial
likelihood does not make an unrestricted hierarchical model more informative
than a flat model.

The proposed model gains statistical strength from its restrictions:
\begin{itemize}[leftmargin=*]
  \item gene abundance is generated by the low-dimensional state \(\bm z_m\);
  \item each \(P_g\) is shared by every cell;
  \item \(C_g\ll T_g\) when a gene has several transcripts;
  \item program weights are generated by \(\bm z_m\) and, optionally, a small
  \(\bm h_m\), rather than estimated independently in every cell.
\end{itemize}
If \(C_g=T_g\), \(P_g\) is the identity, and \(\bm s_{gm}\) is unrestricted,
the model reduces to the flat multinomial parameterization. That limiting
case is not the proposed sparse-cell model.

\section*{Connection to the Paired-Primer Likelihood}

The conceptual transcript abundance remains
\[
  \theta_{tm}
  =
  a_{g(t)m}[P_{g(t)}\bm s_{g(t)m}]_t,
\]
but it need not be instantiated. Project each gene's shared programs through
the primer-specific observation model:
\[
  D_{pkgc}
  =
  \sum_{t\in\mathcal T_g} A_{pkt}P_{gtc}.
\]
The unnormalized expected EC mass is then exactly
\[
  q_{kmp}
  =
  \sum_g\sum_{c=1}^{C_g}
  D_{pkgc}\,a_{gm}s_{gmc}.
\]
This equation is the direct connection between the gene factors, isoform
factors, and observed EC counts. It handles within-gene and cross-gene
ambiguous ECs because all genes represented in row \(k\) contribute to the
sum.

Conditioning on the observed UMI total gives
\[
  \bm n_{mp}\mid N_{mp}
  \sim
  \operatorname{Multinomial}
  \left(
    N_{mp},
    \frac{\bm q_{mp}}{\bm 1'\bm q_{mp}}
  \right).
\]
Ignoring constants, the negative log likelihood is
\[
  \mathcal L
  =
  -\sum_{m,p}\sum_{k:n_{kmp}>0}
    n_{kmp}\log q_{kmp}
  +\sum_{m,p}N_{mp}\log(\bm 1'\bm q_{mp}).
\]
For the normalizer, precompute
\[
  \delta_{pgc}
  =
  \sum_kD_{pkgc}
  =
  \sum_{t\in\mathcal T_g}d_{pt}P_{gtc},
  \qquad
  d_{pt}=\sum_kA_{pkt}.
\]
It follows that
\[
  \bm 1'\bm q_{mp}
  =
  \sum_g\sum_c\delta_{pgc}a_{gm}s_{gmc}.
\]
Neither likelihood term requires a transcript-by-cell abundance block.

The two primer halves share \(\bm z_m\), \(\bm h_m\), \(a_{gm}\), and
\(\bm s_{gm}\), but use separate \(D_p\) obtained from their separate
\(A_p\). Conditioning on \(N_{mp}\) removes a primer-by-cell exposure
parameter and preserves the current convention of normalizing each half
independently.

An independent Poisson model,
\[
  n_{kmp}\sim\operatorname{Poisson}(s_{mp}q_{kmp}),
\]
is an alternative when total UMI information is intended to contribute to
fitting. It requires estimating or profiling the exposure \(s_{mp}\). The
conditional multinomial model is the simpler first implementation and directly
matches the sampling unit used for molecule-count folds.

\section*{Why Cell-Level EM Is Not Required}

Independent EM would estimate a high-dimensional transcript vector in every
cell. That is underdetermined for cells with hundreds or a few thousand UMIs.
The hierarchical model instead estimates only low-dimensional cell
coordinates:
\[
  \bm z_m\in\mathbb R^{r_g},
  \qquad
  \bm h_m\in\mathbb R^{r_i}.
\]
All gene and program parameters are shared across cells.
Cells without molecules that distinguish isoforms remain close to the pooled
program mixture given by \(P_g\) and \(\kappa_{gc}\). Informative cells update
the shared relationship between cell state and program usage.

EM may still be useful on pooled data to initialize the columns of \(P_g\). A
global or pseudobulk EM estimate has enough molecules to establish baseline
isoform usage, but it is not treated as a cell-level estimator. Joint
minibatch gradient optimization of the count likelihood then estimates cell
states, program weights, and shared programs.

\section*{Initialization and Fitting}

A staged fit reduces non-convexity and provides useful diagnostics:
\begin{enumerate}[leftmargin=*]
  \item Aggregate ECs that map entirely within one gene and fit
  \(\alpha_g,\bm u_g,\bm z_m\) with a gene-level count factor model. Cross-gene
  ambiguous ECs are excluded only from this initializer.

  \item Initialize one column of \(P_g\) from pooled EM. Additional columns can
  be initialized from large pseudobulks or small perturbations of the pooled
  profile. Pooled EM is acceptable at this stage.

  \item With \(\bm z_m\) initialized, fit \(\kappa_{gc}\) and \(\bm v_{gc}\)
  under the complete EC likelihood while holding \(r_i=0\).

  \item Jointly fine-tune gene abundance, the program dictionaries \(P_g\),
  program weights, and both primer observation models.

  \item Fit candidate \(r_i>0\) models by adding small isoform-specific factors
  \(\bm h_m,\bm r_{gc}\). Warm-start from smaller to larger \(r_i\).
\end{enumerate}

The fit should use Adam or another stable first-order optimizer initially,
followed by deterministic full-data or large-batch polishing. Gene and isoform
cell factors should use separate learning rates because isoform gradients are
supported by fewer molecules. Gradient clipping and log-sum-exp evaluation
are required for numerical stability.

Weak Gaussian priors or equivalent quadratic penalties on \(\bm h_m\) and
\(\bm r_{gc}\) are appropriate even though the previous direct factorization did
not use an \(L_2\) penalty. Here they control an explicitly optional residual
state that would otherwise fit sparse-cell noise. Their strength, or the
decision to omit \(\bm h_m\), must be selected from held-out molecules rather
than labels.

\section*{Scalable Likelihood Evaluation}

The likelihood constructs neither a dense cell-by-EC matrix nor a
transcript-by-cell abundance block. For a cell minibatch, only ECs with
nonzero counts contribute to \(\sum_k n_{kmp}\log q_{kmp}\). Store
\(D_{pkgc}\) sparsely by EC and gene-program pair. A COO event table containing
cell, primer, EC, and count joins each observed EC to those sparse projected
program entries. Segmented reductions compute
\[
  q_{kmp}=\sum_{g,c}D_{pkgc}a_{gm}s_{gmc}.
\]

The normalizing term uses the much smaller \(\delta_{pgc}\) table:
\[
  \bm 1'\bm q_{mp}
  =
  \sum_{g,c}\delta_{pgc}a_{gm}s_{gmc}.
\]
For a minibatch of \(B\) cells, the main dense workspaces are
\(G\times B\) gene abundances and at most
\((\sum_g C_g)\times B\) program weights. With small \(C_g\), both are smaller
than a \(T\times B\) transcript block.

Because \(D_p\) depends on the learned \(P_g\), it cannot remain permanently
fixed while \(P_g\) is updated. The block-sparse matrix containing all \(P_g\)
has only \(\sum_g T_gC_g\) nonzeros. Sparse multiplication by \(A_p\) can
refresh \(D_p\) after a global-program update, or gradients can flow through
the sparse projection directly. A practical alternating schedule updates cell
states and program weights for several steps with \(D_p\) cached, then updates
\(P_g\) and refreshes \(D_p\).

Explicit cell-factor storage remains small. For example, one million cells
with a 32-dimensional float32 gene state require approximately 128 MB; an
additional eight-dimensional isoform state requires approximately 32 MB.
An amortized encoder from gene counts to \(\bm z_m\) is a later option, but is
not required for the initial million-cell implementation.

\section*{Identifiability}

The following conventions remove avoidable invariances:
\begin{itemize}[leftmargin=*]
  \item Center the gene intercepts \(\alpha_g\), because a common shift of all
  gene log abundances is removed by the conditional multinomial likelihood.

  \item Constrain every column of \(P_g\) to the transcript simplex. Without
  this constraint, scale can move between \(P_g\) and \(a_{gm}\).

  \item Center \(\kappa_{gc}\), each coordinate of \(\bm v_{gc}\), and each
  coordinate of \(\bm r_{gc}\) over programs \(c\) within a gene. Adding a
  common value to every program logit leaves \(\bm s_{gm}\) unchanged.

  \item Resolve program-label permutation for reporting by ordering programs
  by pooled weight and then matching warm-started programs to the preceding
  fit. Permutation does not change the likelihood.

  \item Center and standardize cell-factor columns after each epoch, absorbing
  location and scale changes into intercepts and loadings.

  \item Use factor balancing or an orthogonality convention for reporting.
  Latent rotations do not affect fitted abundance, but a stable convention
  improves warm starts and interpretation.
\end{itemize}

These constraints are parameter-identification conventions, not biological
regularization and not simplex constraints on the genome-wide transcript
matrix.

\section*{Label-Blind Model Selection}

Molecule counts should be partitioned before any normalization, as in the
existing count-fold procedure. Candidate models are compared using held-out
multinomial log likelihood on the same response scale.

The recommended nested path for the independent isoform state is:
\[
  r_i=0,2,4,8,16,
\]
with small-to-large warm starts. The gene rank \(r_g\) can be selected first,
then held fixed while selecting the small program cap \(C_{\max}\), for example
\(C_{\max}=1,2,3\). After selecting \(C_{\max}\), a one-standard-error rule
should prefer \(r_i=0\) or the smallest isoform rank statistically
indistinguishable from the minimum held-out loss. A gene uses
\(C_g\leq\min(C_{\max},T_g)\), and genes without evidence for variable isoform
usage can remain at \(C_g=1\).

Held-out diagnostics should also decompose performance into:
\begin{itemize}[leftmargin=*]
  \item gene-unambiguous molecule likelihood;
  \item within-gene EC likelihood;
  \item cross-gene ambiguous EC likelihood;
  \item poly(dT) and random-hexamer likelihoods.
\end{itemize}
This distinguishes genuine improvement in isoform prediction from improvement
that comes only through gene abundance. Cell-type labels, ARI, NMI, and
silhouette remain external assessments and are never hyperparameter inputs.

\section*{Downstream Differential Isoform Analysis}

The fitted
\(\bm\pi_{gm}=P_g\bm s_{gm}\) is the conditional isoform usage quantity needed
downstream. Differential effects can be introduced in the program logits:
\[
  \lambda_{gmc}
  =
  \kappa_{gc}+\bm v_{gc}'\bm z_m+\bm r_{gc}'\bm h_m
  +\bm x_m'\bm\gamma_{gc},
\]
where \(\bm x_m\) contains diagnosis, cell type, donor covariates, and
interactions. Coefficients \(\bm\gamma_{gc}\) are centered over programs within
each gene. Program-weight contrasts imply transcript-level changes through
\(P_g\), avoiding independent transcript effects that sparse cells cannot
support.
Inference should respect donor replication, either through donor-level random
effects or by fitting and testing on donor-by-cell-type pseudobulks of the
model-implied sufficient statistics.

This preserves the intended division of labor: gene factors provide a robust
cell representation, shared isoform programs borrow information across cells,
and downstream tests operate on transcript-resolved conditional usage rather
than requiring independent cell-level EM estimates.

\section*{Recommended First Implementation}

The first implementation should deliberately omit the free isoform cell state:
\[
  \log a_{gm}=\alpha_g+\bm u_g'\bm z_m,
  \qquad
  \bm\pi_{gm}=P_g\bm s_{gm},
  \qquad
  s_{gmc}
  =
  \operatorname{softmax}_c
  (\kappa_{gc}+\bm v_{gc}'\bm z_m),
\]
with
\[
  q_{kmp}
  =
  \sum_{g,c}D_{pkgc}a_{gm}s_{gmc}.
\]
It should use \(C_{\max}=2\), the paired-primer multinomial EC likelihood,
pooled-EM initialization of one \(P_g\) column, explicit cell factors, and GPU
minibatches. The fit evaluates the likelihood through \(D_p\) and
\(\bm\delta_p\), never through a materialized \(\Theta\). Only after this model
is converged and evaluated should an independent \(\bm h_m\) be added. This
ordering tests the central hypothesis---that robust gene state can support
transcript-resolution inference---without giving every sparse cell enough
freedom to reproduce its EC noise.

\end{document}
